\par\bigskip\noindent\textbf{K101 supporting proof: The separate cubic and Pontryagin-linear component numbers}\par\nobreak
We keep the original \((6,1;2)\), norm-one selected-infinity
comparison of EI, FO, FL, SS, CS and MC, with its actual physical
circle and all original coefficient arrows. This calculation
evaluates the complete fixed-inner pair before the outer-family
and parameter-orientation transfer. Let
\(\kappa=\int_{F_c}c_2(E_c)\) denote the integer of the actual
MC capped positive spinor bundle. This notation names that
existing bundle invariant; its value remains a separate clutch
calculation. No choice of its value enters this ledger.

\par

\medskip\noindent\textbf{FP.1. The two characteristic polynomials}

\par\nobreak

Write \((k_U,\ell_U)\) for the cubic and Pontryagin-linear
integrals of the positive-hyperplane line U, with the same
inner boundary orientation as SS.5 and MC.24. For an isolated
point with the four actual nonzero integral weights \(w_i\),
the weighted link formulas are
\[
 k_U=\frac1{\prod_iw_i},\qquad
 \ell_U=\frac{\sum_iw_i^2}{\prod_iw_i}.                 \tag{FP.1}
\]
For an actual compact fixed curve \(\mathbf P^1\), let the
three normal lines have weights \(w_i\) and degrees \(d_i\)
in their common retained normal-ray order. Expanding the
equivariant Euler denominator gives
\[
 k_U=-\frac{\sum_i d_i/w_i}{\prod_iw_i},\qquad
 \ell_U=\frac{2\sum_iw_id_i-
                   (\sum_iw_i^2)(\sum_i d_i/w_i)}{\prod_iw_i}.
                                                        \tag{FP.2}
\]
Indeed in a formal variable u the normal denominator is
\(\prod_i(w_iu+d_i)\). The coefficient needed from
\(u^3/\prod_i(w_iu+d_i)\) is its term of degree one on
the curve, namely the first expression in FP.2. For
the second expression the numerator is
\(u\sum_i(w_iu+d_i)^2\); the tangent Pontryagin class
of the curve has degree four and restricts to zero.
Multiplying out the degree-one terms gives the displayed
numerator exactly. Thus the finite weight denominators
are retained separately in both integrals.

On a fixed surface with negative and positive normal
lines \(N_-,N_+\), put
\[
 A=c_1(N_-),\quad B=c_1(N_+),\quad
 a_2=\int A^2,\quad a b=\int AB,\quad b_2=\int B^2,
 \quad p=\int p_1(TF).
\]
The positive spinor is \(\overline{N_-}\oplus N_+\),
with \(c_1=B-A\) and \(c_2=-AB\). MC.24\textquotesingle s full
projective-bundle calculation therefore gives
\[
 k_U=-(a_2-a b+b_2),\qquad
 \ell_U=-(p+a_2+2a b+b_2).                             \tag{FP.3}
\]
In each of FP.1-\/-3 the combination
\((4k_U-\ell_U)/24\) agrees with the earlier scoped
normal-link expression, through these explicit
polynomial identities.

\par

\medskip\noindent\textbf{FP.2. The four strict source surfaces and the finite points}

\par\nobreak

The full normal arrows SS.8 give on the orbifold source
\[
 A=-(j+1)L,\qquad B=-(4-j)(O+T)+jL.
\]
Using SS.5\textquotesingle s actual section maps gives
\[
 a_2=0,\qquad a b=\frac{(j+1)(4-j)}2,\qquad
 b_2=\frac{j^2}{2}-8.                                 \tag{FP.4}
\]
At each of the two A3 source points, the negative-line
degree vectors are respectively
\((1,0,0),(1,-1,1),(0,0,1),(0,0,0)\). The positive
line has zero degree on all three exceptional curves.
These are SS.13\textquotesingle s actual monomial normal maps, in their
signed physical order. With the retained integral
intersection matrix \(I=\operatorname{tridiag}(1,-2,1)\),
the negative square changes by \(d^tI^{-1}d\) and the
mixed and positive squares do not change. This is the
relative/absolute intersection comparison of SS.4,
using the same flat boundary character for each line.

At the A1 patch, MF2d-\/-e gives degree one to the odd
normal and degree zero to the other. MF2c identifies
the odd line as negative for even j and positive for
odd j. The corresponding single square changes by
\(-1/2\), using the actual matrix \([-2]\), while
the mixed square remains unchanged. Thus the resolved
triples and tangent numbers are
\[
\begin{array}{c|rrrr|rr}
 j&a_2&a b&b_2&p&k_U&\ell_U\\\hline
 0&-2&2&-8&-24&12&30\\
 1&-2&3&-8&-24&13&28\\
 2&-2&3&-6&-24&11&26\\
 3&0&2&-4&-24&6&24.
\end{array}                                             \tag{FP.5}
\]
Their pair is \((42,108)\). The four isolated points
of the same full mu2 resolution have weights
\((2w,-w,-w,-w)\). FP.1 gives \((-1/2,-7/2)\)
at each, so their pair is \((-2,-14)\). The finite
group identification and all four distinct node
indices remain the ones in MF2; no extra copy is
introduced by this separate evaluation.

\par

\medskip\noindent\textbf{FP.3. The compact O/T component pairs}

\par\nobreak

For each of the fourteen compact surface stars CS.3,
the root certificate uses the fresh auxiliary
cocharacter \((5,13,29,53)\) in the original quotient
lattice. In a chart N its actual vector is
\(N^{-1}(5,13,29,53)^t\). The two zero physical-weight
rows supply the two auxiliary tangent weights
\(t_1,t_2\); the physical weights \(-1,+1\) select
the two ordered auxiliary normal weights a,b.
The full normal coefficient morphisms of CS.13
make these the weights of the actual two line
bundles. Consequently their separate intersection
numbers are
\[
 a_2=\sum\frac{a^2}{t_1t_2},\quad
 a b=\sum\frac{ab}{t_1t_2},\quad
 b_2=\sum\frac{b^2}{t_1t_2},\quad
 p=\sum\frac{t_1^2+t_2^2}{t_1t_2}.                     \tag{FP.6}
\]
This is applied to the complete oriented surface
atlases in CS.14. The circle is defined by the same
torus character on every chart, every tangent weight
is checked to be nonzero, and all actual chart
transition matrices are retained. Thus FP.6 is an
independent evaluation of the full line intersections,
including the mixed term; it does not reconstruct
two separate numbers from their single combination.

For every one of the eighteen interior compact
curves, the same fresh auxiliary circle is read in
its two actual source charts. Their tangent weights
are \(t,-t\). Matching the three normal rays through
the full inverse chart matrices gives endpoint
weights \(n_{i,0},n_{i,\infty}\) and degree
\((n_{i,0}-n_{i,\infty})/t\). This equals the degree
of the literal normal coefficient overlap and is
checked against it. The twelve p-curves and ten
interface curves use their already proved full
clutch degrees CS.6 and FL.6. The points use the
original four quotient-chart physical weights.

The complete pairs for one O or T copy are
\[
\begin{array}{c|rr}
 \text{component block}&k_U&\ell_U\\\hline
 106\text{ interior points and }18\text{ interior curves}&-18&-126\\
 4\text{ exceptional surfaces}&-20&-20\\
 10\text{ refinement surfaces}&18&126\\
 12\text{ p-curves}&0&-18\\
 10\text{ finite-base interface curves}&0&0.
\end{array}                                             \tag{FP.7}
\]
Each component has its separate pair in the certificate,
together with every input normal weight or intersection
number. CS.15\textquotesingle s complete source-period map gives two
copies of precisely these blocks. The four strict
surfaces and four mu2 points in FP.5 are counted once
globally. The minimal-face IDs in the retained atlas
are disjoint between these blocks, and MC.3 supplies
the one remaining moving face.

\par

\medskip\noindent\textbf{FP.4. The moving bundle and the complete fixed-inner identity}

\par\nobreak

Adding FP.5 and the finite points to two copies of
FP.7 gives the exact pair \((0,18)\). MC.21 and
MC.23 give \(c_1(E_c)=0\), \(\int p_1(TF_c)=-6\)
for the actual moving face and full reference cap.
Its separate pair is therefore
\[
 (k_U,\ell_U)_{\rm moving}=(\kappa,4\kappa+6),
 \qquad \kappa=\int_{F_c}c_2(E_c).                     \tag{FP.8}
\]
Here the full U2 clutch remains MC.16a and PL.4-\/-6;
its second Chern number has not been replaced by
an integer selected from a congruence. The complete
fixed-inner identities are consequently
\[
 (k_U,\ell_U)_{\rm inner}=(\kappa,4\kappa+24),
 \qquad \frac{4k_U-\ell_U}{24}=-1.                     \tag{FP.9}
\]
The actual physical line M is dual to U through
the K68 anti-equivariant bundle map, so its inner
pair is the negative of FP.9. The outer boundary
and Jin parameter directions are separate typed
maps, carried in the full outer-transfer proof.
This note does not insert either of those maps
by convention or assign the separate return pair.

The exact program
\path{code/part3_root/e6_selected_infinity_full_pair_ledger.py}
passes 1,027 assertions. It saves the input hashes,
all 296 actual chart inverses, the fresh auxiliary
weights used in every surface and interior curve,
all separate line intersections, every component
pair, the complete block sums and the symbolic
cancellation of the named actual \(\kappa\).
The full geometry of the component domains and
their original boundary arrows is the one proved
in EI, FO, FL, SS, CS and MC; FP.1-\/-9 evaluate
its separate characteristic integrals.

